% This article's macros.
%
% The framework's vocabulary (\lean/\leanok/\uses/\notready, \ledger, \notes, \sscite, and
% the [T]/[A] status tags) lives in linkage-macros.tex, which the `article-kit` repo owns.
% Everything below the rule is this article's own NOTATION — yours to change.
%
% Anything you add here is automatically render-safe (fatal check 4 reads this file's
% \newcommand definitions); anything you use but do not define must be vetted in
% render-allowlist.txt, or the manifest would silently drop it.

% The linkage framework's macro vocabulary. OWNED BY THE `linkage` REPO — do not edit this
% copy; edit it there and re-run `linkage scaffold --sync`. `linkage check` reports drift.
%
% Two tiers live here, and neither is the article's to choose:
%
%   (a) the leanblueprint API (\lean, \leanok, \uses, \notready, \proves), stubbed as no-ops
%       so the blueprint also compiles with plain pdflatex;
%   (b) the linkage macros proper (\ledger, \notes, \sscite, \statedin, \restates, and the
%       [T]/[A] status tags).
%       The checker and the hub manifest key on these names, and the hub reads every
%       satellite's manifest with one reader — so an article that renamed them would emit a
%       manifest nothing could consume.
%
% The article's own NOTATION belongs in its own macros file, alongside an \input of this one.

% -- (a) leanblueprint API ----------------------------------------------------
% When compiled standalone (via print.tex) these are no-ops. When built through the
% plasTeX/leanblueprint pipeline, the package's own definitions take precedence and the
% dependency graph / Lean-link tracking activates. See the load-order note in web.tex.
\providecommand{\lean}[1]{}
\providecommand{\leanok}{}
\providecommand{\uses}[1]{}
\providecommand{\notready}{}
\providecommand{\proves}[1]{}

% -- (b) linkage macros -------------------------------------------------------
% \ledger{A15} : this [A] node is grounded in axiom-ledger entry A15 (renders inline, so it
%                reads the same as the older prose form "ledger A15").
% \notes{slug} : the node's "what" lives in the wiki content note <slug>. A no-op in the PDF,
%                so private wiki slugs stay out of the public build; the checker reads it
%                from the source.
% \sscite{key} : inline citation by pinned citekey (the librarian's Better BibTeX key).
%                Renders as the typewriter @citekey both in the PDF and, via macro expansion
%                at manifest emit, in transcluded wiki blocks — where it matches the hub's
%                claim-ref convention and stays checkable.
% \statedin{line:def:cascade-family} : this node uses a shared-library statement whose home
%                is that node of that article (LINKAGE.md rule 15). A no-op in the PDF, as
%                \lean and \uses are: write the citation a reader sees beside it, as the
%                paper prints it (\cite[Def.~3.2]{...}); the checker resolves this one.
% \restates{line:def:cascade-family} : this node repeats that home's statement so that a
%                reader can follow; the checker compares the two texts. An editorial
%                rendering pins the home's sha: \restates{line:def:cascade-family@<sha12>}.
\newcommand{\ledger}[1]{ledger~#1}
\newcommand{\notes}[1]{}
\newcommand{\sscite}[1]{\texttt{@#1}}
\newcommand{\statedin}[1]{}
\newcommand{\restates}[1]{}

% -- Status tags --------------------------------------------------------------
% [T] proved-target  : Lean must prove this from primitives + the [A] interfaces.
% [A] analytic axiom : stated as a Lean axiom / carried hypothesis; the paper cites the
%     external analysis for it, and the ledger entry records the page anchor.
% Every statement node must declare one (fatal check 5), and every [A] node must say in its
% status annotation what the citation carries and what it does not (fatal check 7).
\newcommand{\Ttag}{\textbf{[T]}\ \textsf{proved-target}}
\newcommand{\Atag}{\textbf{[A]}\ \textsf{analytic interface}}
\newcommand{\statusT}{\hfill{\small$\vartriangleright$ \Ttag}}
\newcommand{\statusA}{\hfill{\small$\vartriangleright$ \Atag}}


% -- Notation -----------------------------------------------------------------
\newcommand{\RR}{\mathbb{R}}
\newcommand{\CC}{\mathbb{C}}
\newcommand{\NN}{\mathbb{N}}
\newcommand{\ZZ}{\mathbb{Z}}
\newcommand{\Cin}{\operatorname{Cin}}   % the entire cosine integral, as in module B
\newcommand{\Levy}{L\'evy}              % the name, accented through one macro
\newcommand{\Matern}{Mat\'ern}          % the same device for the second member's name
% The programme's internal labels of its released papers (Q-0344, the author's review of the
% isotropic paper): \prog{V}{Thm.~7.3} is a numbered place in Paper V's line paper (V) or its cone
% module (B), \progpaper{B} the paper itself. The blueprint prints its short form; each paper defines
% the same two macros to print a real citation (paper*/macros.tex), so a statement shared verbatim
% reads [V, Thm. 7.3] here and [2, Thm. 7.3] in a paper.
\newcommand{\prog}[2]{[#1, #2]}
\newcommand{\progpaper}[1]{[#1]}
% A numbered place of [V] or [B] that the isotropic paper restates in its section 2 (Q-0406, the
% author's review of that paper's section 3: a reader of section 3 is sent to the restatement, which
% carries the released number, not to the earlier paper). \restated{V}{Def.~3.1}{def:line-family}
% prints here exactly what \prog{V}{Def.~3.1} prints; the third argument is the label of the
% restatement in paper-iso/02-line-papers.tex, which a paper's own \restated prints as
% "Definition~\ref{def:line-family}" (and the affine paper as the located citation where it has no
% such label). \restatedp is the same place cited as an aside: [V, Def. 3.1] here, the reference in
% parentheses in a paper. The blueprint has no node with these labels and never \ref's them.
\newcommand{\restated}[3]{[#1, #2]}
\newcommand{\restatedp}[3]{[#1, #2]}
% A clause of such a place, whose number may differ in the restatement:
% \restatedc{B}{Prop.~4.9}{prop:line-filtration}{(4)}{(3)} prints [B, Prop. 4.9(4)] here and
% "Proposition 2.6(3)" in the isotropic paper.
\newcommand{\restatedc}[5]{[#1, #2#4]}
% The affine chapter: parameter sets, the group, and the matrix operations.
\newcommand{\PD}{\mathrm{PD}}           % positive definite d x d matrices
\newcommand{\PSD}{\mathrm{PSD}}         % positive semidefinite d x d matrices
\newcommand{\GL}{\mathrm{GL}}
\newcommand{\Orth}{\mathrm{O}}          % the orthogonal group O(d)
\newcommand{\tp}{\mathsf{T}}            % transpose, written A^{\tp}
\newcommand{\tr}{\operatorname{tr}}
\newcommand{\rank}{\operatorname{rank}}
\newcommand{\diag}{\operatorname{diag}}
\newcommand{\Sym}{\mathrm{Sym}}         % real symmetric d x d matrices
\newcommand{\gl}{\mathfrak{gl}}         % all real d x d matrices, as a Lie algebra
